\section{The Co-Emergence of Matter and Force}

Matter and force are not two separate things bolted onto the mesh. They are two sectors of one dock, woven together by the one knit. This section makes that claim precise and then tests it. The structural half, that the photon sector and the fermion sectors must coexist, follows from the symmetry of the $24$ sites and was settled earlier. The dynamical half, that the sectors actually \emph{act} on each other, is the result developed here. A single coupling ties them. Turn it off and both effects vanish at once.

A word on what is at stake. In ordinary physics the existence of matter, the existence of light, and the coupling between them are three independent inputs. One writes down a fermion field, one writes down a gauge field, and one writes down a term that joins them, each chosen by hand. The mesh does not have that freedom. The same $24$ sites that carry the fermion already carry the photon, and the one knit that evolves them already joins the two. There is no separate dial for matter, for force, and for their union. There is one substrate and one law.

\subsection{The two sectors live in one dock}

The substrate is the $\{3,4,3,4\}$ honeycomb. The $24$ sites out of any dock are the directions of the $24$-cell, equivalently the $D_4$ root system. Under triality the order-three outer symmetry of $\mathrm{SO}(8)$, these $24$ sites split into three eights, as derived where spin and fermions are treated.

\begin{table}[ht]
\centering
\begin{tabular}{@{}lll@{}}
\toprule
sector & size & role \\
\midrule
$8v$ & $8$ & the photon \\
$8s$ & $8$ & one fermion chirality \\
$8c$ & $8$ & the other fermion chirality \\
\bottomrule
\end{tabular}
\caption{The triality split of the $24$ sites of a dock into one vector and two spinor sectors.}
\label{tab:co-emergence-sectors}
\end{table}

The $8v$ is the photon sector. The $8s$ and $8c$ are the two chiralities of the fermion. All three sit in the same $24$ sites. They are not three different docks or three different fields layered over space. They are three faces of one direction set, and triality is the symmetry that relates them.

\begin{definition}[The three sectors of a dock]
The $24$ sites of a dock, taken as the $D_4$ root system, decompose under $\mathrm{SO}(8)$ as $8v + 8s + 8c$. The vector summand $8v$ is the photon. The two spinor summands $8s$ and $8c$ are the two chiralities of the fermion. Triality, the order-three outer symmetry, permutes the three eights.
\end{definition}

\subsection{The structural half, coexistence by symmetry}

That the sectors must coexist is a statement about symmetry, and it was proved first because it needs no dynamics at all.

Two facts settle it. First, every rotation of the sites preserves each sector. A rotation moves directions among themselves, but it does not mix the vector $8v$ with the spinors $8s$ and $8c$. Each eight is closed under the rotation group. Second, triality relates the three eights to one another. The very symmetry that lets us call them by separate names also forbids us from keeping one and discarding the others. They are one orbit under the larger symmetry, not three independent objects.

\begin{proposition}[Coexistence is forced]
Under the rotation symmetry of the dock each sector $8v$, $8s$, $8c$ is preserved. Under triality the three sectors are related as one orbit. Therefore the photon sector and the fermion sectors cannot be separated by hand. Keeping the fermion forces keeping the photon. They co-emerge as a matter of symmetry.
\end{proposition}

This is real and it is necessary, but it has a clear limit. Coexistence by symmetry says the two sectors must be present together. It does not say they \emph{act} on each other. Two fields can share a substrate and still evolve independently, ignoring one another entirely. Whether the photon and the fermion truly interact is a dynamical question, and structure alone cannot answer it.

\subsection{The dynamical half, the sectors act both ways}

The dynamical claim is tested by running one coupled evolution under the one knit and measuring whether the two sectors push on each other. They do, in both directions.

\begin{table}[ht]
\centering
\begin{tabular}{@{}lll@{}}
\toprule
direction & effect & measured signature \\
\midrule
fermion to photon & a moving charge sources a radiating field & field energy grows from exactly zero \\
photon to fermion & a background field accelerates the fermion & acceleration at the Newton rate \\
\bottomrule
\end{tabular}
\caption{The two measured couplings between the sectors, both run under one knit.}
\label{tab:co-emergence-directions}
\end{table}

The first effect is emission. A fermion in motion in the $8s$ or $8c$ sector sources a field in the $8v$ sector. The energy of that radiated field begins at exactly zero and grows as the evolution proceeds. This is a charge radiating a photon field.

The second effect is the Lorentz response. A background field in the $8v$ sector pushes on the fermion. The fermion accelerates, and it does so at the expected Newton rate. This is a field deflecting a charge.

Both effects are measured in one and the same evolution, under one knit, not two laws stitched together. The same fixed table that advances the tones is doing all of the work. Neither effect was inserted as a separate rule for matter or a separate rule for force.

\begin{result}[The sectors interact both ways]
Under one knit a moving fermion sources a radiating photon field whose energy grows from zero, and a background photon field accelerates the fermion at the Newton rate. Both couplings are measured in a single coupled evolution. The photon and the fermion act on each other.
\end{result}

\subsection{The kill switch, one coupling controls both}

The discriminating control is the single coupling constant of the knit. This is the test that could have failed.

If the two effects were artifacts of separate hidden dynamics, each driven by its own machinery, then switching off the one coupling would leave at least one of them standing. The emission might survive, or the deflection might, while the other died. That outcome would mean the sectors were not really bound. They would only have looked bound while both happened to be running.

That is not what happens. At zero coupling the one knit decouples and both effects vanish exactly. No field is sourced. The fermion does not accelerate.

\begin{table}[ht]
\centering
\begin{tabular}{@{}lll@{}}
\toprule
coupling & field sourced & fermion acceleration \\
\midrule
nonzero & yes, energy grows from zero & yes, at the Newton rate \\
zero & none & none \\
\bottomrule
\end{tabular}
\caption{The decoupling control. One coupling ties both effects, and zeroing it kills both together.}
\label{tab:co-emergence-killswitch}
\end{table}

\begin{principle}
At zero coupling the one knit decouples and both effects vanish exactly. The binding is genuine. One constant ties the photon and the fermion, and removing it cleanly separates them.
\end{principle}

Both effects die together because they were never separate. They are two readings of one coupling in one law. This is what distinguishes a real binding from a coincidence of two independent processes. A coincidence could be broken apart one piece at a time. A binding cannot.

\subsection{What the result establishes}

The structural proof showed that the sectors must coexist. It left open whether they interact. This result closes that gap, and it does so on one flat substrate under one knit.

\begin{itemize}
\item The fermion sources the photon.
\item The photon deflects the fermion.
\item One coupling controls both.
\item Zero coupling kills both.
\end{itemize}

\begin{table}[ht]
\centering
\begin{tabular}{@{}ll@{}}
\toprule
item & status \\
\midrule
sectors must coexist by symmetry & proved structurally \\
sectors act on each other both ways & measured, depth L2 \\
decoupling kills both & measured, the control \\
\bottomrule
\end{tabular}
\caption{The depth and status of the co-emergence claim.}
\label{tab:co-emergence-status}
\end{table}

This is the dynamical co-emergence of matter and force. Light and matter are not two ingredients joined by a third. They are two sectors of one dock, woven by one knit, and the single thread that binds them is the same thread that, when cut, sets them free.

One thing this result does not settle is the \emph{magnitude} of the coupling that ties the sectors. The knit fixes the form of the binding and proves that it is real. It does not fix the value of the constant. That value is the analogue of the fine-structure constant, and the honest answer about it is given where the coupling value is treated. The binding is established here. Its size is a separate question, deferred and answered plainly there.
